% ===================================================================
% TEMPLATE de paper acadêmico (agentic-study-latex) — versão probabilística.
% Estrutura: Abstract · Introdução · Metodologia · Resultados · Conclusões · Limitações · Fontes.
% Substitua {{...}} e preencha \tbd{} a partir de dist_stats.md + inventario_fontes.md.
% ANTES de compilar: gere as figuras com monte_carlo_dist.py (dist_itens.png, dist_total.png).
% Compile: latexmk -pdf <arquivo>.tex   |   Pronto quando grep "\tbd{" == 0.
% NAO dependa de siunitx: use \reais{}.
% ===================================================================
\documentclass[11pt,a4paper]{article}
\usepackage[T1]{fontenc}
\usepackage[utf8]{inputenc}
\usepackage[brazilian]{babel}
\usepackage{amsmath, amssymb}
\usepackage{booktabs}
\usepackage{graphicx}
\usepackage{float}
\usepackage[margin=2.5cm]{geometry}
\usepackage{microtype}
\usepackage{xcolor}
\usepackage[hidelinks]{hyperref}
\usepackage{setspace}
\onehalfspacing
\graphicspath{{figuras/}}

\newcommand{\reais}[1]{R\$\,#1}
\newcommand{\tbd}[1]{\textcolor{red}{\textbf{[A PESQUISAR: #1]}}}

\title{\textbf{{{TITULO}}}\\[0.3em]\large {{SUBTITULO}}}
\author{{{AUTOR}}}
\date{\today}

\begin{document}
\maketitle

\begin{abstract}
\noindent {{RESUMO: o que se estima, o método (distribuição de modelos $\times$ preços de mercado,
Monte Carlo), e o resultado: \tbd{E[V] revenda e reposição + banda 90\% + envelope}.}}
\end{abstract}

\section{Introdução}
{{Contexto: o que foi recuperado e por quê medir.}} A resposta não é um número único: estima-se a
\emph{distribuição} de valores plausíveis de cada item. Por \emph{revenda} entende-se a compra do
bem \emph{usado}; por \emph{reposição}, a compra do equivalente \emph{novo} (o custo de repô-lo).

\section{Metodologia}
\subsection{Modelo probabilístico}
O valor de uma unidade da categoria $i$ na base $b\in\{\text{rev},\text{rep}\}$ é a v.a.\ de mistura
\begin{equation}
  V_i^{\,b}\sim\sum_{m\in\mathcal{M}_i} w_{i,m}\,\hat F_{i,m}^{\,b},\qquad
  \mathbb{E}[V_i^{b}]=\sum_m w_{i,m}\,\bar p_{i,m}^{\,b},
  \label{eq:mix}
\end{equation}
com peso $w_{i,m}$ ($\sum_m w_{i,m}=1$) e $\hat F_{i,m}^{\,b}$ a distribuição empírica dos preços
coletados. O total e seus momentos:
\begin{equation}
  V^{\,b}=\sum_i\sum_{k=1}^{N_i} V_{i,k}^{\,b},\quad
  \mathbb{E}[V^{b}]=\sum_i N_i\,\mathbb{E}[V_i^{b}],\quad
  \mathrm{Var}(V^{b})=\sum_i N_i\,\mathrm{Var}(V_i^{b}).
  \label{eq:total}
\end{equation}
Nenhum fator de depreciação arbitrário: a base de revenda já é preço de usado; a de reposição, de novo.

\subsection{Distribuição de modelos por frequência}
{{Pesos $w$ de fonte de FREQUÊNCIA (roubo/vendas/frota) por categoria; declarar ouro (recorte
local) vs proxy.}}

\subsection{Coleta de preços e distribuições empíricas}
{{n observações por modelo (usado e novo); fontes; amostra $\neq$ censo.}}

\subsection{Agregação por Monte Carlo}
Sorteiam-se $N_i$ unidades por categoria (Eq.~\eqref{eq:mix}), somam-se, e repete-se $\sim$50.000
vezes (semente fixa). Reporta-se a média e a banda de 90\%.

\section{Resultados}
\subsection{Distribuição de valor por item}
\begin{figure}[H]\centering
  \includegraphics[width=\linewidth]{dist_itens}
  \caption{Distribuição de probabilidade do valor por unidade, por categoria.}\label{fig:itens}
\end{figure}
\begin{table}[H]\centering\small\begin{tabular}{llrrr}\toprule
Item & base & $\mathbb{E}[V_i]$ & desvio & p5--p95 \\ \midrule
\tbd{item} & revenda & \tbd{} & \tbd{} & \tbd{} \\
           & reposição & \tbd{} & \tbd{} & \tbd{} \\ \bottomrule
\end{tabular}\caption{Momentos da distribuição de valor por unidade (R\$). Fonte: dist\_stats.md.}\end{table}

\subsection{Valor total recuperado}
\begin{figure}[H]\centering
  \includegraphics[width=0.92\linewidth]{dist_total}
  \caption{Distribuição do valor total (Monte Carlo); linha = média, faixa = banda de 90\%.}\label{fig:total}
\end{figure}
\begin{table}[H]\centering\small\begin{tabular}{lrrrr}\toprule
Item & $N_i$ & Contrib.\ revenda & Contrib.\ reposição & \% \\ \midrule
{{item dominante}} & \tbd{} & \tbd{} & \tbd{} & \tbd{} \\ \midrule
\textbf{Total} & — & \tbd{$V^{rev}$} & \tbd{$V^{rep}$} & 100\% \\ \bottomrule
\end{tabular}\caption{Contribuição por categoria ($N_i\cdot\mathbb{E}[V_i]$) e total (R\$).}\end{table}
\noindent A banda de 90\% é incerteza de amostragem; a incerteza estrutural está no envelope
\tbd{[piso, teto]} e na sensibilidade da maior incógnita.

\section{Conclusões}
{{Síntese: intervalo, item dominante, leitura.}}

\section*{Limitações}
\begin{itemize}\setlength\itemsep{0.1em}
  \item {{distribuição é proxy}} \item {{amostras $\neq$ censo}} \item {{caveat específico do caso}}
\end{itemize}

\section*{Fontes}
\small\noindent Inventário (F-T.N) em \texttt{inventario\_fontes.md}; amostras em
\texttt{dados\_distribuicao.json}. {{lista (SOBRENOME/ÓRGÃO. ``Título'', Veículo, Ano)}}
\end{document}
